\setupcolors[state=start]
\setupcolor[hex]
\definecolor[darkblue][h=003366]
\definecolor[darkgray][h=e6e6e6]
\definecolor[redline][h=d9534f]
\definecolor[blueline][h=5bc0de]
\definecolor[greenline][h=5cb85c]

\setuphead[section][style=\bfa, color=darkblue, before=\blank, after=\blank]
\setuphead[subsection][style=\bf, color=black, before={\blank[2*big]}, after=\blank]
\setuphead[subsubsection][style=\bf, color=darkblue, before=\blank, after=\blank]
\setuppapersize[A4][A4]
\setupbodyfont[11pt]
\setuplayout[width=middle, height=middle, topspace=2cm, bottomspace=2cm]

\starttext

\title{\bf Modelos OPP v5: formulación monolítica y formulación en dos etapas}
\blank[big]

%%%%%%%%%%%%%%%%%%%%%%%%%%%%%%%%%%%%%%%%%%%%%%%%%%%%%%%%%%%%%%%%%%%%
\startsection[title={Introducción}]

Ambos modelos resuelven el mismo problema de {\em order-picking}: dado un conjunto de lotes
de pedidos, un grafo de locaciones de almacén y un conjunto de pickers heterogéneos en
velocidad, determinar simultáneamente la ruta de recolección de cada lote, la asignación de
lotes a pickers y la secuenciación temporal, minimizando una función objetivo que combina
tiempos de completitud con penalizaciones de tardanza y adelanto, sujeto a una restricción
de cadena de frío.

Las dos formulaciones difieren en la estrategia de resolución:

\startitemize
\item {\bf OPP v5.1 (monolítico)}: un único MIP que decide ruta, asignación y secuenciación
en conjunto. El tiempo de ruta $r_b$ es una {\em variable} de decisión.
\item {\bf OPP v5.2 (dos etapas)}: la Etapa~1 resuelve el ruteo óptimo de cada lote de
forma independiente (minimizando el tiempo de ruta); la Etapa~2 recibe el tiempo de ruta
como {\em dato} fijo y resuelve asignación, secuenciación y cadena de frío.
\stopitemize

Ambos tienen el mismo valor óptimo: la ruta más corta domina cualquier otra porque el
objetivo no contiene un término de tiempo de ruta y una ruta más corta solo relaja las
restricciones de duración y cadena de frío. Sin embargo, el conjunto de soluciones óptimas
difiere: el modelo monolítico puede devolver rutas no mínimas que alcanzan el mismo
objetivo.

Los archivos fuente de las formulaciones son:
\startitemize
\item \type{models/opp_v5_1_mono/opp_model_v5_1_mono.zpl}
\item \type{models/opp_v5_2_split/opp_model_v5_2_split_stage1_routing.zpl}
\item \type{models/opp_v5_2_split/opp_model_v5_2_split_stage2_scheduling.zpl}
\stopitemize

\stopsection


%%%%%%%%%%%%%%%%%%%%%%%%%%%%%%%%%%%%%%%%%%%%%%%%%%%%%%%%%%%%%%%%%%%%
\startsection[title={Notación común}]

\startsubsection[title={Conjuntos}]
\startformula
\startalign
    \NC \mathcal{N} \NC = \text{Órdenes de clientes} \NR
    \NC \mathcal{V} \NC = \text{Locaciones del almacén, incluidos } v_{start} \text{ (depósito) y } v_{end} \text{ (nodo final)} \NR
    \NC \mathcal{P} \NC = \text{Pickers} \NR
    \NC \mathcal{B} \NC = \text{Lotes factibles de pedidos} \NR
    \NC \mathcal{E} \NC = \text{Arcos válidos del grafo de locaciones } \subseteq \mathcal{V} \times \mathcal{V} \NR
\stopalign
\stopformula
\stopsubsection

\startsubsection[title={Parámetros}]
\startformula
\startalign
    \NC v_{start}, v_{end} \NC = \text{Nodo depósito inicial y nodo final} \NR
    \NC W^{c}, W^{pen}, W^{e}, W^{t} \NC = \text{Pesos de la función objetivo} \NR
    \NC T^{travel} \NC = \text{Tiempo por unidad de distancia (velocidad de referencia)} \NR
    \NC T^{setup} \NC = \text{Tiempo de preparación por lote (independiente de la ruta)} \NR
    \NC T^{pick} \NC = \text{Tiempo de picking por unidad de capacidad (independiente de la ruta)} \NR
    \NC M \NC = \text{Constante Big-M general (dato del escenario)} \NR
    \NC M^{seq} \NC = \text{Big-M para restricciones de secuenciación (calculado por instancia, ver abajo)} \NR
    \NC C_n \NC = \text{Capacidad requerida por la orden } n \in \mathcal{N} \NR
    \NC T^{due}_n \NC = \text{Due date de la orden } n \in \mathcal{N} \NR
    \NC D_{ij} \NC = \text{Distancia entre } i \text{ y } j, \quad (i,j) \in \mathcal{E} \NR
    \NC S_i \NC = \text{Tiempo de servicio en la locación } i \in \mathcal{V} \NR
    \NC V_p \NC = \text{Factor de velocidad del picker } p \in \mathcal{P} \ (1{,}0 = \text{referencia}) \NR
    \NC T^{cold} \NC = \text{Tiempo máximo fuera de cámara de frío} \NR
    \NC \varepsilon \NC = \text{Tolerancia numérica del chequeo de frío } (10^{-9}) \NR
\stopalign
\stopformula

\startformula
\startalign
    \NC A_{ni} \NC = \startmathcases
                    \NC 1 \MC \text{si la orden } n \text{ requiere visitar la locación } i \NR
                    \NC 0 \MC \text{en otro caso} \NR
                 \stopmathcases \NR
    \NC B_{nb} \NC = \startmathcases
                    \NC 1 \MC \text{si la orden } n \text{ pertenece al lote } b \NR
                    \NC 0 \MC \text{en otro caso} \NR
                 \stopmathcases \NR
\stopalign
\stopformula
\stopsubsection

\startsubsection[title={Parámetros derivados}]
\startformula
\startalign
    \NC L_b \NC = \left| \left\{ i \in \mathcal{V} \setminus \{v_{start},v_{end}\} : \sum_{n \in \mathcal{N}} A_{ni} \cdot B_{nb} \ge 1 \right\} \right| \quad \text{(locaciones requeridas por el lote } b \text{)} \NR
    \NC F_b \NC = T^{setup} + T^{pick} \cdot \sum_{n \in \mathcal{N}} C_n \cdot B_{nb} \quad \text{(tiempo fijo del lote, independiente de la ruta y del picker)} \NR
\stopalign
\stopformula

$M^{seq}$ no es una constante del modelo: se lee de los datos, y el armado de los
archivos \type{.dat} lo calcula para cada instancia como el horizonte en que un solo
picker, el más lento, ejecuta todos los lotes seguidos con una ruta de peor caso, redondeado
al entero y con un piso de 500:
\startformula
\startalign
    \NC M^{seq} \NC = \max\left(500,\; \sum_{b \in \mathcal{B}} \left( F_b + \frac{R^{max}}{\min_{p} V_p} \right) + 1 \right) \NR
    \NC R^{max} \NC = T^{travel} \cdot |\mathcal{V}| \cdot \max_{(i,j) \in \mathcal{E}} D_{ij} \;+\; |\mathcal{V}| \cdot \max_{i \in \mathcal{V}} S_i \NR
\stopalign
\stopformula
\stopsubsection

\stopsection


\page
%%%%%%%%%%%%%%%%%%%%%%%%%%%%%%%%%%%%%%%%%%%%%%%%%%%%%%%%%%%%%%%%%%%%
%%%%%%%%%%%%%%%%%%%%%%%%%%%%%%%%%%%%%%%%%%%%%%%%%%%%%%%%%%%%%%%%%%%%
\startsection[title={OPP v5.1 — Formulación monolítica}]

Un único MIP que decide ruta, asignación y secuenciación en conjunto.

\startsubsection[title={Variables}]
\startformula
\startalign
    \NC z_{ijb} \in \{0,1\} \NC \quad \text{1 si el arco } (i,j) \in \mathcal{E} \text{ se usa en la ruta del lote } b \NR
    \NC q_{ijb} \ge 0 \NC \quad \text{Flujo de carga (Gavish-Graves): locaciones requeridas del lote } b \text{ que faltan visitar al recorrer } (i,j) \in \mathcal{E} \NR
    \NC r_b \ge 0 \NC \quad \text{Tiempo de ruta del lote } b \text{ (traslado + servicio; {\em variable})} \NR
    \NC x_{pb} \in \{0,1\} \NC \quad \text{1 si el lote } b \text{ es asignado al picker } p, \quad \forall (p,b) \in \mathcal{P} \times \mathcal{B} \NR
    \NC u_{pb} \ge 0 \NC \quad \text{Tiempo de procesamiento del lote } b \text{ por el picker } p \NR
    \NC \NC \quad \text{con cota superior } u_{pb} \le T^{cold} + \varepsilon \text{ (cadena de frío)} \NR
    \NC c_{pb} \ge 0 \NC \quad \text{Tiempo de completitud del lote } b \text{ en la agenda del picker } p \NR
    \NC c_b \ge 0 \NC \quad \text{Tiempo de completitud del lote } b \text{ (agregado)} \NR
    \NC y_{bk} \in \{0,1\} \NC \quad \text{1 si } b \text{ va antes que } k \text{ en el mismo picker, } b < k \NR
    \NC ta_n, ea_n \ge 0 \NC \quad \text{Tardanza y adelanto de la orden } n \NR
\stopalign
\stopformula
\stopsubsection

\startsubsection[title={Función objetivo}]
\placeformula
\startformula
\min \left( W^{c} \cdot \sum_{b \in \mathcal{B}} c_b + W^{pen} \cdot \left( W^{e} \cdot \sum_{n \in \mathcal{N}} ea_n + W^{t} \cdot \sum_{n \in \mathcal{N}} ta_n \right) \right)
\stopformula
\stopsubsection

\startsubsection[title={Restricciones de ruteo (1)--(12)}]

\startsubsubsection[title={Rutas de lote}]
\startformula
\startalign[n=4, align={right,left,left,left}]
    \NC (1) \NC \sum_{j:(v_{start},j) \in \mathcal{E}} z_{v_{start} j b} = 1 \NC \quad \forall b \in \mathcal{B} \NC \quad \text{(Salida del depósito)} \NR
    \NC (2) \NC \sum_{i:(i,v_{end}) \in \mathcal{E}} z_{i\, v_{end}\, b} = 1 \NC \quad \forall b \in \mathcal{B} \NC \quad \text{(Entrada al nodo final)} \NR
    \NC (3) \NC \sum_{i:(i,v_{start}) \in \mathcal{E}} z_{i\, v_{start}\, b} = 0 \NC \quad \forall b \in \mathcal{B} \NC \quad \text{(Sin retorno al depósito)} \NR
    \NC (4) \NC \sum_{j:(v_{end},j) \in \mathcal{E}} z_{v_{end}\, j\, b} = 0 \NC \quad \forall b \in \mathcal{B} \NC \quad \text{(Sin salida del nodo final)} \NR
\stopalign
\stopformula

Las restricciones (3) y (4) aplican cuando $v_{start} \ne v_{end}$.
\stopsubsubsection

\startsubsubsection[title={Visita y conservación de flujo}]
\startformula
\startalign[n=4, align={right,left,left,left}]
    \NC (5) \NC \sum_{i:(i,j) \in \mathcal{E}} z_{ijb} \ge \min\!\left(1,\, \textstyle\sum_{n} A_{nj} \cdot B_{nb}\right) \NC \quad \forall b,\; j \in \mathcal{V} \setminus \{v_{start},v_{end}\} \NC \quad \text{(Entrada requerida)} \NR
    \NC (6) \NC \sum_{i:(j,i) \in \mathcal{E}} z_{jib} \ge \min\!\left(1,\, \textstyle\sum_{n} A_{nj} \cdot B_{nb}\right) \NC \quad \forall b,\; j \in \mathcal{V} \setminus \{v_{start},v_{end}\} \NC \quad \text{(Salida requerida)} \NR
    \NC (7) \NC \sum_{g:(g,i) \in \mathcal{E}} z_{gib} - \sum_{j:(i,j) \in \mathcal{E}} z_{ijb} = 0 \NC \quad \forall b,\; i \in \mathcal{V} \setminus \{v_{start},v_{end}\} \NC \quad \text{(Conservación)} \NR
    \NC (8) \NC \sum_{j:(j,i) \in \mathcal{E}} z_{jib} \le 1 \NC \quad \forall b,\; i \in \mathcal{V} \setminus \{v_{start}\} \NC \quad \text{(Máx.\ un arco de entrada)} \NR
    \NC (9) \NC \sum_{j:(i,j) \in \mathcal{E}} z_{ijb} \le 1 \NC \quad \forall b,\; i \in \mathcal{V} \setminus \{v_{end}\} \NC \quad \text{(Máx.\ un arco de salida)} \NR
\stopalign
\stopformula
\stopsubsubsection

\startsubsubsection[title={Eliminación de subtours (Gavish-Graves)}]
\startformula
\startalign[n=4, align={right,left,left,left}]
    \NC (10) \NC q_{ijb} \le L_b \cdot z_{ijb} \NC \quad \forall b,\; (i,j) \in \mathcal{E} \NC \quad \text{(Flujo solo en arcos activos)} \NR
    \NC (11) \NC \sum_{j:(v_{start},j) \in \mathcal{E}} q_{v_{start}\, j\, b} = L_b \NC \quad \forall b \in \mathcal{B} \NC \quad \text{(Flujo total desde depósito)} \NR
    \NC (12) \NC \sum_{g:(g,i) \in \mathcal{E}} q_{gib} - \sum_{j:(i,j) \in \mathcal{E}} q_{ijb} = \min\!\left(1,\, \textstyle\sum_{n} A_{ni} \cdot B_{nb}\right) \NC \quad \forall b,\; i \in \mathcal{V} \setminus \{v_{start},v_{end}\} \NC \quad \text{(Balance GG)} \NR
\stopalign
\stopformula
\stopsubsubsection
\stopsubsection

\startsubsection[title={Definición del tiempo de ruta (13)}]
\startformula
\startalign[n=2, align={right,left}]
    \NC (13) \quad r_b = \NC T^{travel} \cdot \displaystyle\sum_{\scriptstyle (i,j) \in \mathcal{E},\; i \ne v_{end},\; j \ne v_{start}} D_{ij} \cdot z_{ijb}
    \;+\; \sum_{i \in \mathcal{V} \setminus \{v_{start},v_{end}\}} S_i \cdot \sum_{j:(i,j) \in \mathcal{E}} z_{ijb} \quad \forall b \in \mathcal{B} \NR
\stopalign
\stopformula
\stopsubsection

\startsubsection[title={Asignación y tiempo de ejecución (14)--(18)}]
\startformula
\startalign[n=4, align={right,left,left,left}]
    \NC (14) \NC \sum_{p \in \mathcal{P}} x_{pb} = 1 \NC \quad \forall b \in \mathcal{B} \NC \quad \text{(Asignación única)} \NR
    \NC (15) \NC F_b + \frac{r_b}{V_p} - M \cdot (1 - x_{pb}) \le u_{pb} \NC \quad \forall p \in \mathcal{P},\; b \in \mathcal{B} \NC \quad \text{(Cota inferior de procesamiento)} \NR
    \NC (16) \NC u_{pb} \le F_b + \frac{r_b}{V_p} + M \cdot (1 - x_{pb}) \NC \quad \forall p \in \mathcal{P},\; b \in \mathcal{B} \NC \quad \text{(Cota superior de procesamiento)} \NR
    \NC (17) \NC u_{pb} \le c_{pb} \NC \quad \forall p \in \mathcal{P},\; b \in \mathcal{B} \NC \quad \text{(Procesamiento precede completitud)} \NR
    \NC (18) \NC c_{pb} - M \cdot (1 - x_{pb}) \le c_b \NC \quad \forall p \in \mathcal{P},\; b \in \mathcal{B} \NC \quad \text{(Completitud del lote)} \NR
\stopalign
\stopformula

La cadena de frío se implementa mediante la cota superior de la variable:
$u_{pb} \le T^{cold} + \varepsilon$. Cuando $x_{pb} = 1$, (15) fuerza
$u_{pb} \ge F_b + r_b / V_p$, por lo que $F_b + r_b / V_p \le T^{cold} + \varepsilon$.
Cuando $x_{pb} = 0$, $u_{pb}$ puede valer 0 y la cota no restringe.
\stopsubsection

\startsubsection[title={Secuenciación sin solapamiento (19)--(20)}]
\startformula
\startalign[n=3, align={right,left,left}]
    \NC (19) \NC c_{pk} \ge c_{pb} + F_k + \frac{r_k}{V_p} - M^{seq} \cdot (3 - x_{pb} - x_{pk} - y_{bk}) \NC \quad \forall p,\; b < k \NR
    \NC (20) \NC c_{pb} \ge c_{pk} + F_b + \frac{r_b}{V_p} - M^{seq} \cdot (2 - x_{pb} - x_{pk} + y_{bk}) \NC \quad \forall p,\; b < k \NR
\stopalign
\stopformula
\stopsubsection

\startsubsection[title={Tardanza y adelanto (21)--(22)}]
\startformula
\startalign[n=4, align={right,left,left,left}]
    \NC (21) \NC c_{pb} - T^{due}_n - M \cdot (2 - x_{pb} - B_{nb}) \le ta_n \NC \quad \forall p,\; b,\; n : B_{nb} = 1 \NC \quad \text{(Tardanza)} \NR
    \NC (22) \NC T^{due}_n - c_{pb} - M \cdot (2 - x_{pb} - B_{nb}) \le ea_n \NC \quad \forall p,\; b,\; n : B_{nb} = 1 \NC \quad \text{(Adelanto)} \NR
\stopalign
\stopformula

Las restricciones se generan solo para pares $(n,b)$ con $B_{nb} = 1$.
\stopsubsection

\stopsection


\page
%%%%%%%%%%%%%%%%%%%%%%%%%%%%%%%%%%%%%%%%%%%%%%%%%%%%%%%%%%%%%%%%%%%%
%%%%%%%%%%%%%%%%%%%%%%%%%%%%%%%%%%%%%%%%%%%%%%%%%%%%%%%%%%%%%%%%%%%%
\startsection[title={OPP v5.2 — Formulación en dos etapas}]

La resolución se descompone en dos MIP secuenciales. La Etapa~1 resuelve el ruteo
independientemente de los pickers; la Etapa~2 recibe los tiempos de ruta y resuelve
asignación, secuenciación y cadena de frío.

%%%%%%%%%%%%%%%%%%%%%%%%%%%%%%%%%%%%%%%%%%%%%%%%%%%%%%%%%%%%%%%%%%%%
\startsubsection[title={Etapa 1: Ruteo por lote}]

Como ningún término de ruteo depende de qué picker ejecuta el lote ni del orden entre
lotes, el ruteo óptimo de cada lote se resuelve de forma independiente.

\startsubsubsection[title={Variables}]
\startformula
\startalign
    \NC z_{ijb} \in \{0,1\} \NC \quad \text{1 si el arco } (i,j) \in \mathcal{E} \text{ se usa en la ruta del lote } b \NR
    \NC q_{ijb} \ge 0 \NC \quad \text{Flujo de carga (Gavish-Graves): locaciones requeridas del lote } b \text{ que faltan visitar al recorrer } (i,j) \in \mathcal{E} \NR
    \NC r_b \ge 0 \NC \quad \text{Tiempo de ruta del lote } b \text{ (traslado + servicio)} \NR
\stopalign
\stopformula
\stopsubsubsection

\startsubsubsection[title={Función objetivo}]
\placeformula
\startformula
\min \sum_{b \in \mathcal{B}} r_b
\stopformula

Como los lotes no interactúan, minimizar la suma equivale a minimizar cada $r_b$ por
separado.
\stopsubsubsection

\startsubsubsection[title={Restricciones (1)--(13)}]
Idénticas a las restricciones (1)--(13) del modelo monolítico. Definen la estructura de
la ruta (camino de $v_{start}$ a $v_{end}$ visitando locaciones requeridas), la eliminación de
subtours por Gavish-Graves y el cálculo de $r_b$.

La Etapa~1 no usa los conjuntos $\mathcal{P}$ ni los parámetros de scheduling
($W^{c}$, $T^{setup}$, etc.). Solo necesita $\mathcal{N}$, $\mathcal{V}$,
$\mathcal{B}$, $\mathcal{E}$, y los parámetros de grafo y lote.
\stopsubsubsection

\startsubsubsection[title={Salida}]
$r_b^*$ para cada $b \in \mathcal{B}$: el tiempo mínimo de traslado + servicio.
Este valor se pasa como dato a la Etapa~2.
\stopsubsubsection
\stopsubsection

%%%%%%%%%%%%%%%%%%%%%%%%%%%%%%%%%%%%%%%%%%%%%%%%%%%%%%%%%%%%%%%%%%%%
\startsubsection[title={Etapa 2: Asignación, secuenciación y cadena de frío}]

Recibe $r_b^*$ como el dato fijo $R_b$. Ya no decide nada sobre el ruteo.

\startsubsubsection[title={Parámetros adicionales (derivados del resultado de Etapa 1)}]
\startformula
\startalign
    \NC R_b \NC = \text{Tiempo de ruta óptimo del lote } b \text{ ({\em dato} de Etapa 1)} \NR
    \NC T^{exec}_{pb} \NC = F_b + \frac{R_b}{V_p} \quad \text{(parámetro precalculado, no variable)} \NR
    \NC \mathcal{PB}^{elig} \NC = \left\{ (p,b) \in \mathcal{P} \times \mathcal{B} : T^{exec}_{pb} \le T^{cold} + \varepsilon \right\} \NR
\stopalign
\stopformula

El conjunto $\mathcal{PB}^{elig}$ implementa la cadena de frío: solo los pares
(picker, lote) cuyo tiempo de ejecución no excede el límite de frío son elegibles. Todas
las variables y restricciones de la Etapa~2 se definen solo sobre $\mathcal{PB}^{elig}$.
\stopsubsubsection

\startsubsubsection[title={Variables}]
\startformula
\startalign
    \NC x_{pb} \in \{0,1\} \NC \quad \text{1 si el lote } b \text{ es asignado al picker } p, \quad (p,b) \in \mathcal{PB}^{elig} \NR
    \NC c_{pb} \ge 0 \NC \quad \text{Tiempo de completitud del lote } b \text{ en el picker } p, \quad (p,b) \in \mathcal{PB}^{elig} \NR
    \NC c_b \ge 0 \NC \quad \text{Tiempo de completitud del lote } b \NR
    \NC y_{bk} \in \{0,1\} \NC \quad \text{1 si } b \text{ va antes que } k \text{ en el mismo picker, } b < k \NR
    \NC ta_n, ea_n \ge 0 \NC \quad \text{Tardanza y adelanto de la orden } n \NR
\stopalign
\stopformula

No hay variable $u_{pb}$: como $T^{exec}_{pb}$ es un parámetro, el vínculo entre tiempo
de ejecución y completitud no requiere Big-M.
\stopsubsubsection

\startsubsubsection[title={Función objetivo}]
\placeformula
\startformula
\min \left( W^{c} \cdot \sum_{b \in \mathcal{B}} c_b + W^{pen} \cdot \left( W^{e} \cdot \sum_{n \in \mathcal{N}} ea_n + W^{t} \cdot \sum_{n \in \mathcal{N}} ta_n \right) \right)
\stopformula

Idéntica a la del modelo monolítico.
\stopsubsubsection

\startsubsubsection[title={Asignación y completitud (1)--(3)}]
\startformula
\startalign[n=4, align={right,left,left,left}]
    \NC (1) \NC \sum_{p: (p,b) \in \mathcal{PB}^{elig}} x_{pb} = 1 \NC \quad \forall b \in \mathcal{B} \NC \quad \text{(Asignación única en pares elegibles)} \NR
    \NC (2) \NC T^{exec}_{pb} \cdot x_{pb} \le c_{pb} \NC \quad \forall (p,b) \in \mathcal{PB}^{elig} \NC \quad \text{(Piso de completitud, sin Big-M)} \NR
    \NC (3) \NC c_{pb} - M \cdot (1 - x_{pb}) \le c_b \NC \quad \forall (p,b) \in \mathcal{PB}^{elig} \NC \quad \text{(Completitud del lote)} \NR
\stopalign
\stopformula

La restricción (2) aprovecha que $T^{exec}_{pb}$ es un coeficiente constante: el producto
$T^{exec}_{pb} \cdot x_{pb}$ es lineal. Cuando $x_{pb} = 1$, impone
$c_{pb} \ge T^{exec}_{pb}$; cuando $x_{pb} = 0$, el lado izquierdo es cero.
Esto elimina la necesidad de la variable $u_{pb}$ y sus dos restricciones Big-M (15)--(16)
del modelo monolítico.
\stopsubsubsection

\startsubsubsection[title={Secuenciación sin solapamiento (4)--(5)}]
\startformula
\startalign[n=3, align={right,left,left}]
    \NC (4) \NC c_{pk} \ge c_{pb} + T^{exec}_{pk} - M^{seq} \cdot (3 - x_{pb} - x_{pk} - y_{bk}) \NC \quad \forall p,\; b < k,\; (p,b),(p,k) \in \mathcal{PB}^{elig} \NR
    \NC (5) \NC c_{pb} \ge c_{pk} + T^{exec}_{pb} - M^{seq} \cdot (2 - x_{pb} - x_{pk} + y_{bk}) \NC \quad \forall p,\; b < k,\; (p,b),(p,k) \in \mathcal{PB}^{elig} \NR
\stopalign
\stopformula

Misma estructura que (19)--(20) del monolítico, usando $T^{exec}_{pb}$ como parámetro
en lugar de la expresión $F_b + r_b / V_p$ (que en el monolítico involucra la variable
$r_b$).
\stopsubsubsection

\startsubsubsection[title={Tardanza y adelanto (6)--(7)}]
\startformula
\startalign[n=4, align={right,left,left,left}]
    \NC (6) \NC c_{pb} - T^{due}_n - M \cdot (2 - x_{pb} - B_{nb}) \le ta_n \NC \quad \forall n,\; (p,b) \in \mathcal{PB}^{elig} : B_{nb} = 1 \NC \quad \text{(Tardanza)} \NR
    \NC (7) \NC T^{due}_n - c_{pb} - M \cdot (2 - x_{pb} - B_{nb}) \le ea_n \NC \quad \forall n,\; (p,b) \in \mathcal{PB}^{elig} : B_{nb} = 1 \NC \quad \text{(Adelanto)} \NR
\stopalign
\stopformula
\stopsubsubsection

\startsubsubsection[title={Dominio de variables (8)--(9)}]
\startformula
\startalign[n=3, align={right,left,left}]
    \NC (8) \NC ta_n,\; ea_n,\; c_b,\; c_{pb} \ge 0 \NC \quad \forall n,\; b,\; (p,b) \in \mathcal{PB}^{elig} \NR
    \NC (9) \NC x_{pb},\; y_{bk} \in \{0,1\} \NC \quad \forall (p,b) \in \mathcal{PB}^{elig},\; b < k \NR
\stopalign
\stopformula
\stopsubsubsection
\stopsubsection

\stopsection


\page
%%%%%%%%%%%%%%%%%%%%%%%%%%%%%%%%%%%%%%%%%%%%%%%%%%%%%%%%%%%%%%%%%%%%
\startsection[title={Diferencias entre las formulaciones}]

\bTABLE[split=yes, align=middle, frame=on, width=broad]
\setupTABLE[row][1][background=color, backgroundcolor=darkgray, style=bold]
\bTR \bTH Aspecto \eTH \bTH v5.1 Monolítico \eTH \bTH v5.2 Dos etapas \eTH \eTR
\bTR \bTD Número de MIPs \eTD \bTD 1 \eTD \bTD 2 (secuenciales) \eTD \eTR
\bTR \bTD Tiempo de ruta ($r_b$ / $R_b$) \eTD \bTD Variable de decisión \eTD \bTD Dato (resultado de Etapa 1) \eTD \eTR
\bTR \bTD $T^{exec}_{pb}$ \eTD \bTD Expresión con variable ($F_b + r_b / V_p$) \eTD \bTD Parámetro precalculado \eTD \eTR
\bTR \bTD Variable $u_{pb}$ \eTD \bTD Sí (con Big-M en (15)--(16)) \eTD \bTD No (innecesaria) \eTD \eTR
\bTR \bTD Cadena de frío \eTD \bTD Cota superior: $u_{pb} \le T^{cold} + \varepsilon$ \eTD \bTD Dominio: $\mathcal{PB}^{elig}$ \eTD \eTR
\bTR \bTD Dominio de $x_{pb}$ \eTD \bTD $\mathcal{P} \times \mathcal{B}$ (todos los pares) \eTD \bTD $\mathcal{PB}^{elig}$ (solo elegibles) \eTD \eTR
\bTR \bTD Piso de completitud \eTD \bTD $u_{pb} \le c_{pb}$ (17) \eTD \bTD $T^{exec}_{pb} \cdot x_{pb} \le c_{pb}$ (2) \eTD \eTR
\bTR \bTD Ruta garantizada mínima \eTD \bTD No (puede elegir ruta no mínima) \eTD \bTD Sí (Etapa 1 minimiza $r_b$) \eTD \eTR
\bTR \bTD Tamaño del MIP principal \eTD \bTD Grande (ruteo + scheduling juntos) \eTD \bTD Menor (cada etapa por separado) \eTD \eTR
\eTABLE

\stopsection

\stoptext
